% ===========================================================================
\chapter{The live map}\label{ch:livemap}

Everything is in place: a map that draws itself at any zoom, three search
algorithms, directions, a navigator, and a traveller to guide. This
chapter puts them all in a window, the way a map app does.

\section{What it does}

The program opens a window showing the whole island. It works in
\emph{stages}:

\begin{enumerate}[nosep]
  \item \textbf{Choosing the start.} Click anywhere on the map, or on a
    place's name. The click snaps to the nearest node of the roads.
  \item \textbf{Choosing the destination.} Click again.
  \item \textbf{Searching.} The search runs, and its explored region
    unfolds on the map over two seconds, in purple, so you can watch it
    spread. Then the route appears in blue.
  \item \textbf{Travelling.} The view zooms in on the traveller and follows
    them as they go, with the next instruction across the top of the
    window. The part of the route already travelled turns grey, and the
    last few GPS fixes show as faint dots around the traveller.
  \item \textbf{Arrived.} Click to choose the next destination; the new
    trip starts where the last one ended.
\end{enumerate}

And keys change things at any time: \code{w} switches between walking and
driving, \code{a} between the three search algorithms, \code{s} between
the shortest and the quickest route, and \code{x} makes the traveller take
a wrong turn. \code{z} switches between following the traveller and
showing the whole map, the space bar pauses, \code{+} and \code{-} speed
the trip up and slow it down, \code{n} starts a new trip, and \code{q}
quits.

\section{An animation loop}

A program with a moving picture works like a film: it draws a still
picture, a \emph{frame}, shows it, and draws the next one, many times a
second. Between frames, it reads the mouse and keyboard and moves
everything on by the time that has passed. In outline:

\begin{shell}
until the window is closed:
    read the mouse and the keyboard
    move everything on by the time since the last frame
    draw the whole frame, off screen
    show it
\end{shell}

The whole frame is drawn every time, map and all: the canvas has no way to
move a picture already drawn. Drawing it takes the canvas around a
hundredth of a second, less when zoomed in, since roads off the canvas are
skipped; with the short pause each frame ends with, that makes something
like forty frames a second.

\emph{Double buffering} is what makes this look smooth. With it on,
drawing goes to an image the window doesn't show, and \code{show} copies
the finished frame into the window in one go. Without it, you would see
each frame being drawn, the map flickering as it is painted over.

\section{The stages and the colours}

\begin{newfile}{navigate.nex}
-- navigate.nex: the live map. Choose where to start and where to go,
-- watch the route being searched for, then watch a traveller follow it
-- while the navigator talks them through every turn.
--
--   nex navigate.nex                        -- click to choose the places
--   nex navigate.nex harbour lighthouse     -- start with these two
--   nex navigate.nex harbour lighthouse --walk --algorithm dijkstra
--
-- Keys, while the window is open:
--   w        walk or drive              a        change search algorithm
--   s        shortest or quickest way   x        take a wrong turn
--   z        follow the traveller or show the whole map
--   space    pause                      + / -    faster / slower
--   n        new trip                   q        quit
--
-- Without a display, it draws the start of the trip into navigate.png.
-- For pictures of a trip without opening a window:
--   --snapshot SECONDS FILE   draw the trip SECONDS into the journey
--   --wrong-turn SECONDS      and make the traveller go wrong then
--   --search-snapshot FILE    draw the search, partway through

intern Simulation
intern Map_View
intern graphics/Key
intern time/Date_Time

enum union Stage
  Choosing_Start
  Choosing_Finish
  Searching
  Travelling
  Arrived
end

class Colours
feature
  ROUTE = create Color.make(40, 110, 220)
  ROUTE_CASING = create Color.make(20, 60, 140)
  TRAVELLED = create Color.make(150, 160, 175)
  EXPLORED = create Color.make_rgba(150, 60, 200, 160)
  START = create Color.make(40, 160, 70)
  FINISH = create Color.make(210, 50, 50)
  TRAVELLER = create Color.make(20, 90, 200)
  FIX = create Color.make_rgba(20, 90, 200, 90)
  PANEL = create Color.make_rgba(255, 255, 255, 230)
  WARNING = create Color.make_rgba(255, 235, 200, 240)
  TEXT = create Color.make(40, 40, 40)
  DIM = create Color.make(110, 110, 110)
end
\end{newfile}

\section{The program's state}

Everything the program knows lives in one object, of class \code{App}: the
map and its view, a graph for each mode of travel, the current settings,
the stage, the places chosen, the latest search, the trip, and a few things
the animation needs.

\begin{append}{navigate.nex}
class App
create
  make(m: Island_Map) do
    island := m
    view := create Map_View.make(m, 1100, 760)
    view.set_insets(64.0, 26.0)
    view.canvas.set_title(m.name)
    driving := create Road_Graph.build(m, Mode.Driving)
    walking := create Road_Graph.build(m, Mode.Walking)
    graph := driving
    algorithm := "astar"
    shortest := false
    stage := Stage.Choosing_Start
    start_node := 0
    finish_node := 0
    start_name := ""
    finish_name := ""
    revealed := 0
    hold := 0.0
    speedup := 1.0
    following := true
    paused := false
    quitting := false
    was_pressed := false
    alert := 0.0
    seen_reroutes := 0
    trail := []
  end
feature
  island: Island_Map
  view: Map_View
  driving: Road_Graph
  walking: Road_Graph
  -- The graph for the current way of travelling: driving or walking.
  graph: Road_Graph
  algorithm: String
  shortest: Boolean
  stage: Stage
  start_node: Integer
  finish_node: Integer
  start_name: String
  finish_name: String
  -- The latest search, and how many of its explored nodes are on screen.
  search: ?Search_Result
  revealed: Integer
  -- Seconds to wait, showing the route, before the traveller sets off.
  hold: Real
  trip: ?Trip
  -- How many seconds of the trip pass in one second on screen.
  speedup: Real
  -- Whether the view follows the traveller, zoomed in.
  following: Boolean
  paused: Boolean
  quitting: Boolean
  was_pressed: Boolean
  -- Seconds left to show the "off route" warning.
  alert: Real
  seen_reroutes: Integer
  -- The latest GPS fixes, oldest first.
  trail: Array[Point]

  -- How long the search takes to unfold on screen, and how long a whole
  -- trip takes at the starting speed, in seconds.
  SEARCH_SECONDS = 2.0
  TRIP_SECONDS = 30.0
  -- The zoom while following the traveller, in metres per pixel.
  CLOSE_UP = 1.4
\end{append}

The panels across the top and bottom of the window are 64 and 26 pixels
high, and \code{set_insets} tells the view, so that the map is centred in
the space between them. The window is titled with the map's name. The search and the trip are \code{?} types,
because there aren't any until the places have been chosen.

A real trip takes minutes; nobody wants to watch one in real time. So the
trip runs faster than real time, by \code{speedup}, which is chosen for
each trip so that it takes about 30 seconds on screen.

\section{Planning a trip}

A trip can start from two place names, given on the command line, or from
two clicks. Either way, \code{plan} runs the search, prints the directions
in the terminal, and moves to the searching stage. Changing the mode of
travel can leave the start or the destination off the new graph (a
footpath has no node in the driving graph), so \code{plan} snaps both to
the nearest node of the current graph first.

\begin{continued}{navigate.nex}
  use_algorithm(name: String) do
    algorithm := name
  end

  -- Start a trip between two places, given by (part of) their names.
  choose_places(from_name, to_name: String) do
    if ?island.find_place(from_name) as a and ?island.find_place(to_name) as b then
      start_name := a.name
      start_node := graph.nearest_node(island.point(a.node_id))
      finish_name := b.name
      finish_node := graph.nearest_node(island.point(b.node_id))
      plan()
    else
      raise "unknown place: " + from_name + " or " + to_name
    end
  end

  -- Search for a route from the start to the finish, and show the search
  -- unfolding.
  plan() do
    -- After a change of mode, the start and finish move to the nearest
    -- nodes of the new graph.
    start_node := graph.nearest_node(island.point(start_node))
    finish_node := graph.nearest_node(island.point(finish_node))
    search := find_route(graph, start_node, finish_node, algorithm, shortest)
    revealed := 0
    trip := nil
    trail := []
    stage := Stage.Searching
    if ?search as s then
      print_directions(s)
    end
  end
\end{continued}

\section{The mouse and the keyboard}

The canvas tells us whether a mouse button is down right now, but not
whether it has just been clicked. A click is the moment the button goes
down, so we remember whether it was down at the last frame, and act only
when it is down now and wasn't then.

Typed keys queue up between frames, and we take them all.

\begin{continued}{navigate.nex}
  -- Read the mouse and the keyboard.
  handle_input() do
    let c: Canvas := view.canvas
    let pressed: Boolean := c.is_mouse_pressed()
    if pressed and not was_pressed and c.has_mouse_position() then
      clicked(create Point.make(c.mouse_x(), c.mouse_y()))
    end
    was_pressed := pressed
    from
    until
      not c.has_next_key_typed()
    do
      key_typed(c.next_key_typed())
    end
    if c.is_key_pressed(Key.ESCAPE) then
      quitting := true
    end
  end
\end{continued}

\code{mouse_x} and \code{mouse_y} give the mouse's position in the
canvas's own coordinates, which are our metres. So a click is already a
point on the map, ready to snap to the nearest node. If it is near a
place's dot, it takes the place's name; otherwise it is ``the marked spot''.

What a click means depends on the stage. The first sets the start, the
second the destination, which starts the search. After arriving, a click
sets a new destination, starting from where the traveller is. Clicking the
start again just moves it.

\begin{continued}{navigate.nex}
  clicked(p: Point) do
    let name: String := name_near(p)
    let node: Integer := graph.nearest_node(p)
    if stage = Stage.Choosing_Start then
      start_node := node
      start_name := name
      stage := Stage.Choosing_Finish
    elseif stage = Stage.Choosing_Finish or stage = Stage.Arrived then
      if stage = Stage.Arrived then
        start_node := finish_node
        start_name := finish_name
      end
      if node = start_node then
        start_node := node
        start_name := name
        stage := Stage.Choosing_Finish
      else
        finish_node := node
        finish_name := name
        plan()
      end
    end
  end
\end{continued}

Keys change a setting, and the three that change the route (mode,
algorithm, and shortest or quickest) plan it again at once, if there is a
route to plan. \code{graph == driving} asks whether \code{graph} is the
very same object as \code{driving}, which is how we tell which mode we are
in.

\begin{continued}{navigate.nex}
  key_typed(k: String) do
    if k = "q" then
      quitting := true
    elseif k = " " then
      paused := not paused
    elseif k = "+" or k = "=" then
      speedup := speedup * 2.0
    elseif k = "-" then
      speedup := (speedup / 2.0).max(1.0)
    elseif k = "z" then
      following := not following
    elseif k = "n" then
      stage := Stage.Choosing_Start
      search := nil
      trip := nil
    elseif k = "x" then
      if ?trip as t then
        t.traveller.take_wrong_turn()
      end
    elseif k = "w" or k = "a" or k = "s" then
      if k = "w" then
        graph := when graph == driving then walking else driving end
      elseif k = "a" then
        if algorithm = "astar" then
          algorithm := "bfs"
        elseif algorithm = "bfs" then
          algorithm := "dijkstra"
        else
          algorithm := "astar"
        end
      else
        shortest := not shortest
      end
      if stage /= Stage.Choosing_Start and stage /= Stage.Choosing_Finish then
        plan()
      end
    end
  end
\end{continued}

\section{Moving on}

\code{update} moves everything on by the real time since the last frame.

While searching, it reveals more of the search's explored nodes: enough
each second that all of them are on screen after \code{SEARCH_SECONDS}.
Then the trip begins, after a pause of a second and a half to let you see
the route. While travelling, it moves the trip on by the elapsed time
multiplied by the speed-up, and notes each new GPS fix.

\begin{continued}{navigate.nex}
  -- Move everything on by `seconds` of real time.
  update(seconds: Real) do
    alert := (alert - seconds).max(0.0)
    if not paused and ?search as s then
      if stage = Stage.Searching then
        let per_second: Real := s.explored.length.to_real / SEARCH_SECONDS
        revealed := (revealed + (per_second * seconds).ceiling.max(1)).min(s.explored.length)
        if revealed = s.explored.length then
          if s.path.length < 2 then
            stage := Stage.Arrived
          else
            stage := Stage.Travelling
            hold := 1.5
            let t: Trip := create Trip.make(graph, s.path, finish_name, algorithm, 42)
            speedup := (t.navigator.route.duration() / TRIP_SECONDS).max(1.0)
            trip := t
            seen_reroutes := 0
          end
        end
      elseif stage = Stage.Travelling and ?trip as t then
        if hold > 0.0 then
          hold := hold - seconds
        else
          t.advance(seconds * speedup)
          note_fix(t)
          if t.is_over() then
            stage := Stage.Arrived
          end
        end
      end
    end
  end

  -- Remember the trip's latest GPS fix, and whether the navigator has just
  -- planned a new route.
  note_fix(t: Trip) do
    if trail.is_empty or trail.get(trail.length - 1).distance_to(t.fix) > 0.0 then
      trail.add(t.fix)
      if trail.length > 40 then
        trail := trail.drop(1)
      end
    end
    if t.navigator.reroutes > seen_reroutes then
      seen_reroutes := t.navigator.reroutes
      alert := 2.5
      console.print_line("Off the route: planning a new one.")
    end
  end
\end{continued}

When the navigator plans a new route, the program shows a warning in the
top panel for two and a half seconds, by setting \code{alert}, which
\code{update} counts down.

\section{Drawing a frame}

A frame is drawn in layers, like the map itself: the map, then the search
or the route, then the markers for the start and destination, then the
traveller, then the panels on top. While travelling with \code{following}
on, the view first moves to centre on the traveller, zoomed in; at every
other time, it shows the whole island.

\begin{continued}{navigate.nex}
  -- Draw one frame.
  draw() do
    if following and stage = Stage.Travelling and ?trip as t then
      view.look_at(t.traveller.position(), CLOSE_UP)
    else
      view.show_all()
    end
    view.draw_map()
    if stage = Stage.Searching then
      draw_explored()
    end
    if stage = Stage.Travelling or stage = Stage.Arrived then
      draw_route()
    end
    draw_markers()
    if ?trip as t then
      draw_traveller(t)
    end
    draw_banner()
    draw_status()
  end
\end{continued}

The search is drawn as in the try-its: a line from each explored node to
its parent, and a dot on each node, for as many explored nodes as have been
revealed so far. When all of them have, the route goes on top.

\begin{continued}{navigate.nex}
  -- The search so far: a line from each explored node to the node the
  -- search reached it from.
  draw_explored() do
    if ?search as s then
      let c: Canvas := view.canvas
      c.set_pen_color(Colours.EXPLORED)
      c.set_pen_radius(2.5)
      across s.explored.take(revealed) as n do
        if s.parents.contains_key(n) then
          let a: Point := island.point(s.parents.get(n))
          let b: Point := island.point(n)
          c.line(a.x, a.y, b.x, b.y)
        end
      end
      c.set_pen_radius(1.0)
      across s.explored.take(revealed) as n do
        let p: Point := island.point(n)
        c.filled_circle(p.x, p.y, view.metres(3.0))
      end
      if revealed = s.explored.length and s.found() then
        view.line(island.points_of(s.path), Colours.ROUTE, 5.0)
      end
    end
  end

  -- The route: the part already travelled in grey, the rest in blue.
  draw_route() do
    if ?trip as t then
      let r: Route_Line := t.navigator.route
      let d: Real := t.navigator.progress
      view.line(r.part(d, false), Colours.TRAVELLED, 6.0)
      view.line(r.part(d, true), Colours.ROUTE_CASING, 9.0)
      view.line(r.part(d, true), Colours.ROUTE, 6.0)
    elseif ?search as s then
      view.line(island.points_of(s.path), Colours.ROUTE, 6.0)
    end
  end
\end{continued}

The markers are map pins: a circle with a letter, A for the start and B for
the destination, on a triangle whose point is on the node itself. The
traveller is a blue dot with a white arrow showing which way they are
heading, and around it the recent GPS fixes, as faint dots.

\begin{continued}{navigate.nex}
  draw_markers() do
    if stage /= Stage.Choosing_Start then
      pin(island.point(start_node), Colours.START, "A")
    end
    if stage /= Stage.Choosing_Start and stage /= Stage.Choosing_Finish then
      pin(island.point(finish_node), Colours.FINISH, "B")
    end
  end

  -- A map pin with its point at p.
  pin(p: Point, colour: Color, letter: String) do
    let c: Canvas := view.canvas
    let r: Real := view.metres(10.0)
    let head: Point := p + create Point.make(0.0, view.metres(22.0))
    c.set_pen_color(colour)
    c.filled_polygon([p.x, head.x - r * 0.8, head.x + r * 0.8],
      [p.y, head.y - r * 0.5, head.y - r * 0.5])
    c.filled_circle(head.x, head.y, r)
    c.set_pen_color(Color.WHITE)
    c.set_font(Font.SANS.bold().with_size(12))
    c.text(head.x, head.y, letter)
  end

  -- The traveller as a blue dot with an arrow, and the recent GPS fixes
  -- as faint dots around their path.
  draw_traveller(t: Trip) do
    let c: Canvas := view.canvas
    c.set_pen_color(Colours.FIX)
    across trail as f do
      c.filled_circle(f.x, f.y, view.metres(3.0))
    end
    let p: Point := t.traveller.position()
    c.set_pen_color(Color.WHITE)
    c.filled_circle(p.x, p.y, view.metres(10.0))
    c.set_pen_color(Colours.TRAVELLER)
    c.filled_circle(p.x, p.y, view.metres(7.5))
    c.set_pen_color(Color.WHITE)
    view.arrowhead(p + direction(t.traveller.heading()) * view.metres(5.0),
      t.traveller.heading(), 8.0)
  end
\end{continued}

\section{The panels}

The panel across the top says what to do now. What that is depends on the
stage: an invitation to click, the progress of the search, the navigator's
prompt with an arrow showing the turn, or the news that you have arrived.
Below it, smaller, a detail line: the distance and time to go, for
instance.

\begin{continued}{navigate.nex}
  -- The panel along the top: what to do next.
  draw_banner() do
    let c: Canvas := view.canvas
    let top: Real := c.y_max
    let left: Real := c.x_min
    let height: Real := view.metres(64.0)
    let width: Real := c.x_max - c.x_min
    let headline: String := ""
    let detail: String := ""
    let turn: Real := 0.0
    let warn: Boolean := alert > 0.0
    if stage = Stage.Choosing_Start then
      headline := "Click where you want to start"
      detail := "or on one of the places"
    elseif stage = Stage.Choosing_Finish then
      headline := "Now click where you want to go"
      detail := "Starting at " + start_name
    elseif stage = Stage.Searching then
      headline := search_name() + " is searching..."
      detail := revealed.to_string + " of " + graph.nodes.length.to_string + " nodes explored"
    elseif stage = Stage.Travelling and ?trip as t then
      let nav: Navigator := t.navigator
      if warn then
        headline := "Off the route: planning a new one"
      elseif hold > 0.0 then
        headline := nav.instructions.get(0).text
      else
        headline := nav.prompt()
        turn := nav.next_instruction().turn
      end
      detail := distance_text(nav.remaining_distance()) + " to go, about "
        + duration_text(nav.remaining_time())
      if nav.reroutes > 0 then
        detail := detail + "   (new route " + nav.reroutes.to_string + "x)"
      end
    else
      headline := "You have arrived at " + finish_name
      detail := "Click to choose where to go next"
    end
    c.set_pen_color(when warn then Colours.WARNING else Colours.PANEL end)
    c.filled_rectangle(left + width / 2.0, top - height / 2.0, width / 2.0, height / 2.0)
    let text_left: Real := left + view.metres(16.0)
    if stage = Stage.Travelling and not warn and hold <= 0.0 then
      draw_turn_arrow(create Point.make(left + view.metres(36.0), top - height / 2.0), turn)
      text_left := left + view.metres(70.0)
    end
    c.set_pen_color(Colours.TEXT)
    c.set_font(Font.SANS.bold().with_size(20))
    c.text_left(text_left, top - view.metres(22.0), headline)
    c.set_pen_color(Colours.DIM)
    c.set_font(Font.SANS.with_size(13))
    c.text_left(text_left, top - view.metres(47.0), detail)
  end
\end{continued}

Everything in the panels is placed relative to the canvas's edges,
\code{x_min}, \code{y_max} and so on, with sizes in pixels converted by
\code{metres}, so the panels stay put however the map below them pans and
zooms.

The turn arrow is the one on every satnav: a stem coming up from below,
bending by the turn's angle, with an arrowhead on the end.
\code{direction(angle)} points along the turn: straight up for 0, to the
right for a right turn.

\begin{continued}{navigate.nex}
  -- An arrow showing a turn of `angle` degrees: up for straight on,
  -- bending right or left.
  draw_turn_arrow(at: Point, angle: Real) do
    let c: Canvas := view.canvas
    let stem: Point := at - create Point.make(0.0, view.metres(20.0))
    let bend: Point := at
    let tip: Point := bend + direction(angle) * view.metres(16.0)
    view.line([stem, bend, tip], Colours.ROUTE, 6.0)
    c.set_pen_color(Colours.ROUTE)
    view.arrowhead(tip + direction(angle) * view.metres(8.0), angle, 12.0)
  end
\end{continued}

The bottom panel shows the settings on the left and the keys on the right.

\begin{continued}{navigate.nex}
  -- The line along the bottom: the settings, and the keys.
  draw_status() do
    let c: Canvas := view.canvas
    let height: Real := view.metres(26.0)
    let width: Real := c.x_max - c.x_min
    let bottom: Real := c.y_min
    c.set_pen_color(Colours.PANEL)
    c.filled_rectangle(c.x_min + width / 2.0, bottom + height / 2.0, width / 2.0, height / 2.0)
    let how: String := when graph == driving then "Driving" else "Walking" end
    let what: String := when shortest then "shortest" else "quickest" end
    let status: String := how + ", " + what + " way, " + search_name()
    if ?search as s then
      status := status + " explored " + s.explored.length.to_string + " nodes"
    end
    if stage = Stage.Travelling then
      status := status + ", " + speedup.format(0) + "x speed"
      if paused then
        status := status + " (paused)"
      end
    end
    c.set_pen_color(Colours.TEXT)
    c.set_font(Font.SANS.bold().with_size(12))
    c.text_left(c.x_min + view.metres(12.0), bottom + height / 2.0, status)
    c.set_pen_color(Colours.DIM)
    c.set_font(Font.SANS.with_size(12))
    c.text_right(c.x_max - view.metres(12.0), bottom + height / 2.0,
      "w walk/drive  a algorithm  s shortest  x wrong turn  z zoom  space pause  +/- speed  n new  q quit")
  end
\end{continued}

Three small helpers finish the class: the algorithm's name for people to
read, the name of the place nearest a click (within 25 pixels), and the
directions printed in the terminal whenever a route is planned.

\begin{continued}{navigate.nex}
  search_name(): String do
    if algorithm = "bfs" then
      result := "Breadth-first search"
    elseif algorithm = "dijkstra" then
      result := "Dijkstra"
    else
      result := "A*"
    end
  end

  -- The name of the place at p, if p is on or near one.
  name_near(p: Point): String do
    result := "the marked spot"
    let best: Real := view.metres(25.0)
    across island.places as place do
      let d: Real := island.point(place.node_id).distance_to(p)
      if d < best then
        best := d
        result := place.name
      end
    end
  end

  print_directions(s: Search_Result) do
    console.print_line("")
    if s.found() then
      console.print_line(start_name + " to " + finish_name + ": "
        + distance_text(route_cost(graph, s.path, by_distance())) + ", about "
        + duration_text(route_cost(graph, s.path, by_time())) + " ("
        + search_name() + " explored " + s.explored.length.to_string + " nodes)")
      across directions(graph, s.path, finish_name) as step do
        console.print_line("  " + step.text)
      end
    else
      console.print_line("No way from " + start_name + " to " + finish_name)
    end
  end
end
\end{continued}

\section{The loop}

Now the animation loop itself. Each time round, it measures the real time
since the last frame, with the clock from \code{time/Date_Time}, and moves
everything on by that much. It never moves on by more than a tenth of a
second at once: if the computer stalls for a moment, the trip pauses with
it, instead of jumping ahead. A short \code{pause} at the end of each frame
gives the rest of the computer a turn.

\begin{append}{navigate.nex}
-- Run the window until it is closed.
function run_window(app: App) do
  let c: Canvas := app.view.canvas
  c.enable_double_buffering()
  app.draw()
  app.view.show()
  let last: Integer := (create Date_Time.now()).millis
  from
  until
    app.quitting or not c.is_window_open()
  do
    app.handle_input()
    let now: Integer := (create Date_Time.now()).millis
    app.update(((now - last).to_real / 1000.0).min(0.1))
    last := now
    app.draw()
    app.view.show()
    c.pause(15)
  end
end
\end{append}

The first \code{show} opens the window; the loop runs until the window is
closed or the user quits.

Not every computer has a display: a server, or a build machine like the
one that made this tutorial's pictures. For those, and for pictures, the
program can also run the app without a window, moving it on in steps of its
own and saving a frame. One function draws the search partway through; the
other draws the trip a given number of seconds into the journey, with a
wrong turn on the way if asked for.

\begin{continued}{navigate.nex}
-- Draw the search partway through, without a window.
function search_snapshot(app: App, file: String) do
  app.update(App.SEARCH_SECONDS * 0.6)
  app.draw()
  app.view.save(file)
  console.print_line("saved " + file)
end

-- Draw the trip as it is `seconds` into the journey, without a window.
function snapshot(app: App, seconds, wrong_turn: Real, file: String) do
  app.update(App.SEARCH_SECONDS + 1.0)
  app.update(2.0)
  if ?app.trip as t then
    let step: Real := 0.25
    let clock: Real := 0.0
    from
    until
      clock >= seconds or t.is_over()
    do
      if wrong_turn >= 0.0 and clock <= wrong_turn and clock + step > wrong_turn then
        t.traveller.take_wrong_turn()
      end
      t.advance(step)
      app.note_fix(t)
      clock := clock + step
    end
    if t.is_over() then
      app.update(0.0)
    end
  end
  app.draw()
  app.view.save(file)
  console.print_line("saved " + file)
end
\end{continued}

In \code{snapshot}, the first \code{update} finishes the search and starts
the trip, and the second waits out the pause before the traveller sets off.
After that, the trip is moved on directly, a quarter of a second at a time.

Last, \code{main} reads the command line, makes the app, and either opens
the window or takes a picture.

\begin{continued}{navigate.nex}
function main() do
  let args: Array[String] := program.arguments
  let names: Array[String] := []
  let walk: Boolean := false
  let algorithm: String := "astar"
  let shot_seconds: Real := -1.0
  let shot_file: String := ""
  let search_file: String := ""
  let wrong_turn: Real := -1.0
  from
    let i: Integer := 0
  until
    i >= args.length
  do
    let arg: String := args.get(i)
    if arg = "--walk" then
      walk := true
    elseif arg = "--algorithm" and i + 1 < args.length then
      algorithm := args.get(i + 1)
      i := i + 1
    elseif arg = "--snapshot" and i + 2 < args.length then
      shot_seconds := args.get(i + 1).to_real
      shot_file := args.get(i + 2)
      i := i + 2
    elseif arg = "--search-snapshot" and i + 1 < args.length then
      search_file := args.get(i + 1)
      i := i + 1
    elseif arg = "--wrong-turn" and i + 1 < args.length then
      wrong_turn := args.get(i + 1).to_real
      i := i + 1
    else
      names.add(arg)
    end
    i := i + 1
  end
  let app: App := create App.make(create Island_Map.load(user_path("kestrel_island.map")))
  if walk then
    app.key_typed("w")
  end
  if algorithm = "bfs" or algorithm = "dijkstra" or algorithm = "astar" then
    app.use_algorithm(algorithm)
  end
  if names.length = 2 then
    app.choose_places(names.get(0), names.get(1))
  end
  if search_file /= "" then
    search_snapshot(app, search_file)
  elseif shot_file /= "" then
    snapshot(app, shot_seconds, wrong_turn, shot_file)
  elseif app.view.canvas.has_display() then
    run_window(app)
  else
    console.print_line("No display: drawing navigate.png instead of opening a window.")
    snapshot(app, 0.0, -1.0, "navigate.png")
  end
end

main()
\end{continued}

\section{Off you go}

Run it:

\begin{shell}
nex navigate.nex
\end{shell}

and the island appears in a window, waiting for two clicks. Or name the two
places, and the search starts at once:

\begin{shell}
nex navigate.nex "beach car park" lighthouse
\end{shell}

Here is that search partway through, as the program drew it. From the car
park, A* has gone up the west coast on the Coast Road, which is heading for
the lighthouse, and also a little way up the lanes past the lake and into
town, which look promising too, at first:

\runoutput{navigate.nex "beach car park" lighthouse --search-snapshot nav_search.png}{nav_search}
\mappicture{nav_search}

The terminal shows the directions as soon as the route is planned. Then the
view zooms in, and the traveller sets off. Here is a trip from Kestrel
School to Old Mill, 25 seconds in. The traveller has come up Anchor Street
and turned onto the Coast Road, and the next turn is onto Hill Road:

\runoutput{navigate.nex school mill --snapshot 25 nav_travel.png}{nav_travel}
\mappicture{nav_travel}

The grey line behind the traveller is the way they have come, and the
blue one ahead is the way still to go; the faint dots are the last few GPS
fixes, scattered a few metres either side of the road. The panel at the
top counts down to the next turn, with its arrow.

Now press \code{x}. Here is the trip from the harbour to the lighthouse,
moments after a wrong turn:

\runoutput{navigate.nex harbour lighthouse --snapshot 55 nav_reroute.png --wrong-turn 40}{nav_reroute}
\mappicture{nav_reroute}

The traveller has turned off Harbour Street into Church Street. The fixes
trail off along it, the navigator has noticed, and the new route leads up
Net Lane and back.

Try the other keys too. Press \code{a} to watch breadth-first search, then
Dijkstra's algorithm, then A* again, on the same trip: the difference in
how much of the island each explores is plain to see. Press \code{w} to
walk, and see the route change to the footpaths, the trip take ten times
as long, and the speed-up rise to match. Press \code{s} for the shortest
route instead of the quickest. And press \code{z} to watch the whole trip
from above.
